\section{Why fixing the initial state to zero is not a free simplification}
\label{sec:zero-start-obstruction}

An initialized iterator carries three pieces of data: its counter, initial
state, and update. A tempting deletion is to fix the initial state to zero.
The following standard-model obstruction explains why the resulting language
cannot simply inherit the current term compiler or selector-based logical
bootstrap. It treats a different operation from the positive iterator
$h^{n+1}(a)$ studied in the reviewed \texttt{checker\_and\_barriers} and
\texttt{prevalidated\_elimination} packages: here zero iterations are allowed,
but the seed is always zero. This observation is new relative to those
reviewed analyses; no literature-wide priority or new iPRA presentation is
claimed.

\paragraph{Two precise candidate languages.}
In the \emph{scalar} language the numerical terms are
\[
 t ::= x\mid 0\mid S(t)\mid R(n;u.h),
 \qquad
 R(0;u.h)=0,\qquad R(Sn;u.h)=h[R(n;u.h)/u].
 \tag{ZS}
\]
Only the update $h$ binds $u$; the counter $n$ does not. In the
\emph{counter-aware} variant the last constructor is $R(n;k,u.h)$ and its
equations are
\[
 R(0;k,u.h)=0,\qquad
 R(Sn;k,u.h)=h[n/k,R(n;k,u.h)/u].
 \tag{ZC}
\]
Both $k$ and $u$ are bound only in $h$. Every term is finite staged syntax;
there is no function-variable sort or uniform evaluator. All assertions
below concern the indicated standard interpretation in $\mathbb N$, with
$S(n)=n+1$. Passive parameters are held fixed during a run. No equality
test, conditional operator, additional arithmetic function, or conversion
from formulas into numerical values belongs to either grammar.

\paragraph{Monotonicity theorem.}
\emph{Every term of either language denotes a coordinatewise nondecreasing
function of all its free numerical variables.}

\paragraph{Proof.}
Use external structural induction on the finite syntax. Variables, zero,
and successor satisfy the assertion. For the recursion clause, the
induction hypothesis makes the counter and the update nondecreasing in
all their arguments.

For a fixed parameter assignment $\rho$, write $a_j^\rho$ for the state
after $j$ updates. In the counter-aware case,
\[
 a_0^\rho=0,\qquad
 a_{j+1}^\rho=h(j,a_j^\rho,\rho).
\]
The state sequence is nondecreasing: $a_0^\rho\le a_1^\rho$ because
zero is least, and from $a_j^\rho\le a_{j+1}^\rho$ monotonicity of
$h$ in both its counter and state gives
\[
 a_{j+1}^\rho=h(j,a_j^\rho,\rho)
 \le h(j+1,a_{j+1}^\rho,\rho)=a_{j+2}^\rho.
\]
For parameter assignments $\rho\le\sigma$, induction over the common
update index gives $a_j^\rho\le a_j^\sigma$. The counters satisfy
$n(\rho)\le n(\sigma)$, so
\[
 a_{n(\rho)}^\rho
 \le a_{n(\rho)}^\sigma
 \le a_{n(\sigma)}^\sigma.
\]
This proves the recursion clause. The scalar case is the same argument
with the update's counter argument omitted. It also covers a free
variable occurring both in the counter expression and as a passive
parameter of the update.\hfill$\square$

\paragraph{Exact zero-profile theorem for the scalar language.}
For a scalar term $t(x_1,\ldots,x_r)$, put
$Z(t)=\{\vec x:t(\vec x)=0\}$.
\emph{Its zero set is either the whole domain or a union of coordinate
zero-hyperplanes:}
\[
 Z(t)=\mathbb N^r
 \quad\hbox{or}\quad
 Z(t)=\bigcup_{j\in J}\{\vec x:x_j=0\}
 \qquad(J\subseteq\{1,\ldots,r\}).
 \tag{Zero profile}
\]
The union may be empty. A profile can be computed by structural recursion
without evaluating the represented function on unbounded inputs.

\paragraph{Proof.}
Zero has the whole-domain profile, a variable has its one coordinate
hyperplane, and a successor term has the empty profile. For the remaining
constructor, fix the passive parameters and let $b=h(0)$. If $b=0$,
every iterate from zero is zero. If $b>0$, monotonicity of $h$ gives
$h(v)\ge h(0)=b>0$ at every natural state $v$; consequently every
positive iterate is nonzero. Thus, pointwise,
\[
 R(n;u.h)=0
 \quad\Longleftrightarrow\quad
 n=0\ \lor\ h[0/u]=0.
 \tag{Zero recurrence}
\]
Hence $Z(R(n;u.h))=Z(n)\cup Z(h[0/u])$. The update is a proper syntax
subterm. Its previously computed profile, with $u$ set to zero, becomes
the whole domain if it includes the $u=0$ hyperplane, and otherwise keeps
the same passive-coordinate hyperplanes. Union with the counter's
profile preserves the displayed classification.

Conversely, every displayed kind of zero set is realized by a scalar
term. Zero and $S0$ realize the whole and empty sets. The constant-update
gate $G(c,a)=R(c;u.a)$ satisfies
$Z(G(c,a))=Z(c)\cup Z(a)$, so nested gates realize any prescribed finite
union of coordinate hyperplanes. This converse classifies zero sets,
not all numerical functions sharing a zero set.\hfill$\square$

\paragraph{Consequences for the proposed deletion.}
In the scalar language no term denotes predecessor: the unary function
$\max(n-1,0)$ has zero set $\{0,1\}$, whereas a unary term has zero set
only $\varnothing$, $\{0\}$, or $\mathbb N$. Nor does a binary term
denote ordinary addition, whose zero set is the singleton $(0,0)$;
that set is neither a union of the two coordinate hyperplanes, nor empty,
nor the whole domain.

In \emph{both} languages the function equal to one at zero and zero at
every positive input is unavailable, because it is decreasing from zero
to one. The initialized selector
\[
 C(0,a,b)=a,\qquad C(Sn,a,b)=b
\]
is therefore unavailable too: specializing $a=1,b=0$ would give that
zero-test. In particular, there is no term compiler that recovers every
initialized iterator from either proposed zero-start grammar on the
same numerical domain with the standard readings of zero and successor.
This conclusion already follows from a constant-update instance, before
questions about pairing or general primitive recursion arise.

The extra old-counter argument makes a real difference. In the
counter-aware language
\[
 P(n):=R(n;k,u.k)
 \quad\text{satisfies}\quad P(0)=0,\qquad P(Sn)=n.
\]
Thus predecessor is a term there, and the scalar zero-profile theorem
does not extend to that language. Its proof fails exactly because the
update can change with the counter: the first update can be zero while
a later update is positive. The monotonicity theorem still applies, and
still rules out the selector and the zero-test. No claim about
term-definability of addition in the counter-aware language is needed
or established here.

\paragraph{What this does not prove.}
These are standard-model, same-domain \emph{term-definability}
obstructions. They do not rule out function graphs using quantified
formulas, definable subdomains, tuple domains, quotient equality, or
some other bi-interpretation permitted in the question. Such a proposal
would need its own proof of totality and functionality of the interpreted
operations, transfer of the restricted induction schema, closure under
newly defined updates, and definable isomorphisms for both composites.
The argument above establishes none of those positive requirements and
is not presented as an impossibility theorem for all simpler
axiomatizations of iPRA.
